\uuid{KWX9}
\titre{Noyau/Image/Théorème du rang/Polynôme}
\auteur{hanani ; mzari}
\organisation{exo7}
\licence{CC BY-NC-SA 4.0}
\datecreate{2020-01-20}
\niveau{L1}
\module{Algèbre}
\chapitre{Application linéaire}
\sousChapitre{Définition}
\difficulte{3}
\theme{Noyau et image, Noyau, Image, Théorème du rang, Polynôme}
\exo7qcmid{lille-2.applications-lineaires.021}

\contenu{
\begin{question}
\qclassification{107.01}
\qtags{motcle=Noyau/Image/Théorème du rang/Polynôme}
On note $\Rr_n[X]$ l'espace des polynômes à coefficients réels de degré $\le n$, $n\in \Nn$. On considère l'application linéaire :
$$\begin{array}{rccc}f:&\Rr_3[X]&\to&\Rr_2[X]\\
& P&\to & R,  \end{array}$$
où $R$ est le reste de la division euclidienne de $P$ par $(X+1)^3$. Quelles sont les assertions vraies ?
\begin{answers}
\bad{$\{X^3\}$ est une base de $\ker f $}.
\good{$\dim \ker f =1$}.
\bad{$\{1+X+X^2\}$ est une base de $\Im f $}.
\bad{$\mbox{rg} (f) =3$}.
\end{answers}
\begin{explanations}  $\ker f=\{P \in \Rr_3[X] \; ; \; R=0\}=\{(X+1)^3Q \; ; \; Q\in\Rr\}$. Donc $\{(1+X)^3\}$ est une base de $\ker f$. D'après le théorème du rang, $\mbox{rg}(f)=\dim \Im f=3$ et l'on vérifie que $\{1,X,X^2\}$ est une base de $\Im f$.
\end{explanations}
\end{question}
}
